\documentclass[10pt]{mypackage}

\usepackage[uprightgreeks,light,noDcommand]{kpfonts}
\renewcommand{\coloneq}{\coloneqq}
\renewcommand{\eqcolon}{\eqqcolon}
%\renewcommand{\mathbb}{\mathds}
%\usepackage{homework}
\usepackage{notes}

\usepackage[ backend=bibtex, style = alphabetic, sorting=ynt ]{biblatex}
\addbibresource{all_references.bib}

\usepackage{parskip}

\fancyhf{}
\fancyhead[R]{Avinash Iyer}
\fancyhead[L]{Countable Borel Relations and Cost}
\fancyfoot[C]{\thepage}

\setcounter{secnumdepth}{0}

\begin{document}
\RaggedRight
\section{Countable Borel Relations, Invariant Measures, and Graphings}%
\subsection{Equivalence Relation Preliminaries}%
First, we collect some definitions.
\begin{definition}\hfill
  \begin{itemize}
    \item A relation $R$ on a set $X$ is a set of ordered pairs $R\subseteq X^2$. We will write $xRy$ to denote $\left( x,y \right)\in R$. We also define the \textit{sections} 
      \begin{align*}
        R_x &= \set{y : \left( x,y \right)\in R}\\
        R^y &= \set{x : \left( x,y \right)\in R}
      \end{align*}
      and the inverse relation
      \begin{align*}
        R^{-1} &= \set{\left( y,x \right) : \left( x,y \right)\in R}.
      \end{align*}
    \item If $X$ is equipped with a Borel structure, then the relation $R$ (i.e., equivalence relations, partial orders, total orders, etc.) is called Borel if $R\subseteq X^2$ is Borel.
    \item We will let $\dom\left( f \right)$, $\ran\left( f \right)$, and $\graph\left( f \right)$ denote the domain, range, and graph of $f$ respectively.
    \item A graph $G$ with vertex set $X$ is a non-reflexive, symmetric relation on $X$. The neighbors of $x\in X$ in the graph $G$ are the set of all $y\in X$ such that $\left( x,y \right)\in G$. The cardinality of the set of neighbors of $X$ is called the degree of $x$, denoted $d_G(x)$.
    \item A path in $G$ is a finite sequence of vertices $\left( x_0,x_1,\dots, x_n \right)$ where $\left( x_i,x_{i+1} \right)\in G$, and $x_i\neq x_j$ whenever $i\neq j$ except for possibly $x_n = x_0$. The equivalence classes of the relation $xEy$ if and only if there exists a path from $x$ to $y$ defines the connected components of the graph. If a graph has one connected component, we say the graph is connected.
    \item If $x_n = x_0$ with $n\geq 3$, then we call the path a cycle. A graph that has no cycles is called acyclic, and a connected acyclic graph is called a tree. A rooted tree is a tree with a distinguished vertex $x_0$ called its root.
    \item If $X$ is a standard Borel space and $E$ is a Borel equivalence relation on $X$, then we say the relation is countable if $\left[ x \right]_E$ is countable for each $x\in X$.
    \item If $\Gamma\curvearrowright X$ is a countable group acting on $X$ such that $\Gamma$ preserves the Borel structure of $X$, then we define the orbit equivalence relation to be $xE_{\Gamma}^X y$ if and only if there exists $g$ such that $g\cdot x = y$.
    \item The orbit space $X/\Gamma$ is defined to be the quotient $X/E_{\Gamma}^X$.
    \item If $E$ is a countable Borel relation, then the \textit{full group} of $E$, denoted $\left[ E \right]$, is the set of Borel automorphisms whose graphs are contained in $E$. The set of partial Borel automorphisms is denoted $\left[ \left[ E \right] \right]$; that is,
      \begin{align*}
        \left[ \left[ E \right] \right] &= \set{\phi\colon A\rightarrow B : \phi\text{ is a bijection, }A,B\subseteq X\text{ are Borel, }\graph\left( \phi \right)\subseteq E}.
      \end{align*}
    \item We call a countable equivalence relation $E$ on $X$ \textit{aperiodic} if every equivalence class $\left[ x \right]_E$ is infinite.
    \item Given Borel equivalence relations $E$ and $F$ on $X$ and $Y$ respectively, we say $E$ is \textit{Borel reducible} to $F$, denoted $E\leq_{B}F$, if there is a Borel map $f\colon X\rightarrow Y$ such that $xEy$ if and only if $f(x)Ff(y)$.
  \end{itemize}
\end{definition}
\begin{example}
  If $X^{\Gamma}$ denote the space of functions of $\Gamma$ onto $X$, then the left shift action of $\Gamma$ on $X^{\Gamma}$ is defined by $\gamma\cdot p\left( \delta \right) = p\left( \gamma^{-1}\delta \right)$.

  If $G$ is a standard Borel group and $\Gamma\subseteq G$ is a countable discrete subgroup (also known as a lattice), then the action of $\Gamma$ onto $G$ by left translation, taking $\gamma\cdot g = \gamma g$, yields the quotient space $G/\Gamma$ given by the set of \textit{right} cosets of $\Gamma$.
\end{example}
\begin{example}
  The following are examples of countable Borel equivalence relations.
  \begin{enumerate}[(i)]
    \item If $X = 2^{\N}$, then the \textit{eventual equality} relation says that $xE_0 y$ if and only if there exists $m$ such that for all $n\geq m$, we have $x_n = y_n$. The tail equivalence relation says that $xE_t y$ if and only if there exist $m$ and $k$ such that for all $n$, we have $x_{m+n} = y_{k + n}$.
    \item If $X = \R$, then the \textit{Vitali relation} says that $xE_v y$ if and only if $x-y\in \Q$.
    \item If $X = \R_{+}$, then the \textit{commensurability relation} says that $xE_c y$ if and only if $x/y\in \Q$.
    \item If $k\geq 2$ and $X$ is the space of subshifts of $k^{\Z}$, where a subshift is a nonempty closed subset of $k^{\Z}$ invariant under the shift map $\left( S(x) \right)_{i} = x_{i-1}$.

      This is a compact subspace of the hyperspace of compact subsets of $k^{\Z}$, so it is a compact and metrizable space. Define two subshifts to be isomorphic if there is a homeomorphism between the two closed subsets that commutes with the shift, and define an equivalence relation as such.
  \end{enumerate}
\end{example}
\begin{theorem}[Feldman--Moore]
  Let $E$ be a countable Borel equivalence relation on $X$. Then, there is a countable group $\Gamma$ and a Borel action of $\Gamma$ on $X$ such that $E_{\Gamma}^X = E$. 
\end{theorem}
\begin{proposition}[Marker Lemma]
  Let $E$ be an aperiodic countable Borel equivalence relation on $X$. Then, there is a sequence $\set{S_n}_{n\geq 1}$ of Borel sets $S_n\subseteq X$ such that
  \begin{enumerate}[(i)]
    \item $S_0\supseteq S_1\supseteq S_2\supseteq\cdots$;
    \item $\bigcap_{n\geq 1}S_n = \emptyset$;
    \item every $S_n$ is a complete section for $E$ (i.e., it has nonempty intersection with each equivalence class of $E$).
  \end{enumerate}
  Such a sequence is called a \textit{vanishing sequence of markers}.
\end{proposition}
\begin{proof}
  Via some results from descriptive set theory, we may assume that the equivalence relation $E$ lives on the set $X = 2^{\N}$. Then, we let $s_n(x)$ be the lexicographically smallest $s\in 2^{n}$ such that $\left\vert [x]_E\cap \mathcal{N}_s \right\vert= \infty$, where $\mathcal{N}_s$ denotes the cylinder set in $X$ corresponding to $s$.

  Define $x\in A_n$ if and only if $x|_{n} = s_n(x)$. Then, we have that $\set{A_n}_{n\geq 1}$ is a decreasing sequence of Borel sets which intersects each class of $E$ at infinitely many points. Then, $A = \bigcap_{n\in \N}A_n$ intersects each class of $E$ in at most one point, so that $B_n = A_n\setminus A$ is the vanishing sequence of markers as desired.
\end{proof}
\subsection{Invariant Measures}%
\begin{definition}
  If $E$ is a countable Borel equivalence relation on $X$, then we say $\mu$ is $E$-invariant if for all $f\in \left[ E \right]$, we have $f_{\ast}\mu = \mu$, where $f_{\ast}\mu$ is the pushforward measure, defined by
  \begin{align*}
    f_{\ast}\mu(A) &= \mu\left( f^{-1}\left( A \right) \right).
  \end{align*}
\end{definition}
\begin{proposition}
  The following are equivalent:
  \begin{enumerate}[(a)]
    \item $\mu$ is $E$-invariant;
    \item $\mu$ is $\Gamma$-invariant for any countable group $\Gamma$ such that $E = E_{\Gamma}^X$;
    \item $\mu$ is $\Gamma$-invariant for some $\Gamma$ such that $E = E_{\Gamma}^X$;
    \item for all $\phi\in \left[ \left[ E \right] \right]$, we have $\mu\left( \dom(\phi) \right) = \mu\left( \ran\left( \phi \right) \right)$.
  \end{enumerate}
\end{proposition}
If $\mu$ is invariant under $E$, then we may define a measure $ \overline{\mu} $ on $E$ by taking
\begin{align*}
  \overline{\mu}\left( A \right) &= \int_{}^{} \left\vert A_x \right\vert\:d\mu(x),
\end{align*}
where again,
\begin{align*}
  A_x &= \set{y: \left( x,y \right)\in A}.
\end{align*}
It is also possible to define
\begin{align*}
  \overline{\mu}'\left( A \right) &= \int_{}^{} \left\vert A^y \right\vert\:d\mu(y),
\end{align*}
but since $\mu$ is invariant, $ \overline{\mu} = \overline{\mu}' $. This follows from the fact that
\begin{align*}
  E &= \bigcup_{i\in \N} \graph\left( g_i \right)
\end{align*}
for some collection of relation-preserving automorphisms $g_i\in \left[ E \right]$. If $A\subseteq E$ is any Borel subset, then
\begin{align*}
  A &= \bigcup_{i\in \N} \left( \graph\left( g_i \right)\cap A \right),
\end{align*}
so there are $f_i\in \left[ \left[ E \right] \right]$ such that $\graph\left( f_i \right) = \graph\left( g_i \right)\cap A$. Therefore, we may write any $A\subseteq E$ as a countable disjoint union of graphs of functions $f\in \left[ \left[ E \right] \right]$. Therefore, it suffices to show that $ \overline{\mu}\left( \graph(f) \right) = \overline{\mu}'\left( \graph\left( f \right) \right)$ for any $f\in \left[ \left[ E \right] \right]$. If $C = \dom(f)$ and $D = \ran(f)$, then $ \overline{\mu}\left( \graph(f) \right) = \mu\left( C \right) $ and $ \overline{\mu}'\left( \graph\left( f \right) \right) = \mu\left( D \right) $, so by the definition of $\left[ \left[ E \right] \right]$, we are done.

Now, we may define an equivalence relation $ \widetilde{E} $ on $E$ by taking $\left( x,y \right)E\left( z,w \right)$ if and only if $xEz$, or equivalently, $xEyEzEw$.
\begin{proposition}
  The invariant measure $M$ on $E$ is $\widetilde{E}$-invariant.
\end{proposition}
\begin{proof}
  Let $\Gamma$ act on $X$ such that $E = E_{\Gamma}^X$. Then, if we define the action $\left( g,h \right)\cdot \left( x,y \right) = \left( g\cdot x,h\cdot y \right)$, then we have that $ E^{E}_{\Gamma^2} = \widetilde{E} $, so it suffices to show that the action preserves $M$.

  For this, if $A\subseteq E$ is Borel, we have
  \begin{align*}
    M\left( \left( g,h \right)\cdot A \right) &= \int_{}^{} \left\vert \set{\left( g\cdot x,h\cdot y \right) : \left( x,y \right)\in A} \right\vert\:d\mu(x).
  \end{align*}
  Setting $B = \left( g,h \right)\cdot A$, then we have $y\in B_x$ if and only if $\left( x,y \right)\in B$, which holds if and only if $\left( g^{-1}\cdot x,h^{-1}\cdot y \right)\in A$, which holds if and only if $h^{-1}\cdot y\in A_{g^{-1}\cdot x}$. Therefore, $B_x = h\cdot A_{g^{-1}\cdot x}$, so that
  \begin{align*}
    M\left( \left( g,h \right)\cdot A \right) &= \int_{}^{} \left\vert A_{g^{-1}\cdot x} \right\vert\:d\mu(x)\\
                                              &= \int_{}^{} \left\vert A_x \right\vert\:d\mu(x)\\
                                              &= M\left( A \right),
  \end{align*}
  where we use the fact that $\mu$ is invariant under $\Gamma$.
\end{proof}
\subsection{Graphings and Cost}%
\begin{definition}
  If $E\subseteq X^2$ is a countable Borel equivalence relation, then a (Borel) \textit{graphing} of $E$ is a graph $ G $ such that the connected components of $G$ are the equivalence classes of $E$. We say that $G$ generates $E$.
\end{definition}
\begin{definition}
  A Levitt graph (L-graph) is a countable family $ \left[ \left[ E \right] \right]\supseteq \Phi=\set{\varphi_i}_{i\in I}$ of partial Borel isomorphisms $\varphi_i\colon A_i\rightarrow B_i$, where $A_i,B_i$ are Borel subsets of $X$.

  We call $\Phi$ \textit{finite} if $I$ is finite. We say $\Phi$ is an Levitt graphing (L-graphing) if $\Phi$ generates $E$ in the sense that $xEy$ if and only if $x = y$ or there is a finite set $i_1,\dots,i_k\in I$ and $\ve_1,\dots,\ve_k\in \set{1,-1}$ such that $x = \varphi_{i_1}^{\ve_1}\cdots \varphi_{i_k}^{\ve_k}(y)$.
\end{definition}
For every L-graph $\Phi$, there is an associated graph $G_{\Phi}$ that generates the same equivalence relation. We say $xG_{\Phi}y$ if and only if $x\neq y$ and there exists $i$ such that $\varphi_i(x) = y$ or $\varphi_i(y) = x$.

Conversely, for every graph $G$, there is an L-graph $\Phi_G$ such that $G = G_{\Phi_G}$. This follows from fixing a Borel (total) ordering $<$ on $X$, and finding a family of partial Borel functions $\set{f_i}_{i\in \N}$ with disjoint graphs such that for all $x\in X$, we have
\begin{align*}
  \set{\left( x,y \right) : xGy\text{ and }x < y} &= \bigcup_{i\in \N}\graph\left( f_i \right).
\end{align*}
Since the $f_i$ are countable-to-one, we have $\set{g_{i,j}}_{j\in K_i}$ a family of partial Borel automorphisms with disjoint graphs such that
\begin{align*}
  \graph\left( f_i \right)^{-1}&= \bigcup_{j\in K_i} \graph\left( g_{i,j} \right.
\end{align*}
(Recall that we define the inverse of a relation to be the coordinates flipped.) Define $\Phi_G = \set{g_{i,j}}_{i,j}$.

Considering the case of an $E$-invariant $\mu$, if $G$ is a graphing of $E|_{A}$ for some conull $E$-invariant Borel subset $A$, then we say $G$ is a graphing of $E$ almost everywhere. Similarly, we may define an L-graphing almost everywhere.

Now, we will discuss the theory of cost.
\begin{definition}
  Let $E$ be a countable Borel equivalence relation on $X$, and let $\mu$ be an $E$-invariant measure with corresponding measure on $E$ given by $M$. If $G\subseteq E$ is a subset, then we define its \textit{cost} to be
  \begin{align*}
    C_{\mu} \left( G \right) &= \frac{1}{2}M\left( G \right).
  \end{align*}
\end{definition}
Observe that $0\leq C_{\mu}\left( G \right) \leq \infty$, and $0 < C_{\mu}\left( G \right)$ if $G$ is a graphing of $E$ provided that it's not the case that almost every $E$-class is a singleton. This follows from the fact that $G$ is a complete section for the relation $ \widetilde{E} $ on $E$, meaning that $M(G) > 0$.

Now, if $\Phi = \set{\varphi_i}_{i\in I}$ is an L-graph, define its cost by taking
\begin{align*}
  C_{\mu}\left( \Phi \right) = \sum_{i\in I}\mu\left( \dom\left( \varphi_i \right) \right) = \sum_{i\in I}\mu\left( \ran\left( \varphi_i \right) \right).
\end{align*}
Then, we have
\begin{align*}
  C_{\mu}\left( \Phi \right) &= \frac{1}{2} \int_{}^{} \sum_{i\in I}\left( \chi_{A_i}(x) + \chi_{B_i}(x) \right)\:d\mu(x),
\end{align*}
where we have $A_i = \dom\left( \varphi_i \right)$ and $B_i = \ran\left( \varphi_i \right)$. 
\nocite{topics_in_orbit_equivalence,classical_descriptive_set_theory,theory_countable_borel_equivalence_relations}
\printbibliography
\end{document}
